\documentclass[UTF8]{ctexart}
\usepackage{graphicx}
\usepackage[a4paper]{geometry}
\geometry{left=2.0cm,right=2.0cm,top=2.0cm,bottom=2.0cm}
\usepackage{fancyhdr}
\usepackage{ulem}
\usepackage{amssymb}
\usepackage{color}
\usepackage{xcolor}
\usepackage{mdframed}
\usepackage{mhchem}
\usepackage{siunitx}
\usepackage{booktabs}

\definecolor{myblue}{RGB}{0,0,139}
\definecolor{myyellow}{RGB}{254,244,209}
\newcommand\bluebf[1]{\heiti \textcolor{myblue}{\textbf{#1}} \rmfamily}

\makeatletter
\newcommand\dlmu[2][6cm]{\hskip1pt\underline{\hb@xt@ #1{\hss#2\hss}}\hskip3pt}
\makeatother

%实验信息
\def\name{铂电极表面的电化学反应}
\def\mytoday{\number\year .\number\month .\number\day}
\def\expday{2019年9月18日}
\def\exptmp{26.1$^{\circ}$C}
\def\expprs{102.54 kPa}
\def\zhaiyao{本实验通过循环伏安法，在不同扫描速度下绘制了被氮气饱和的硫酸溶液中铂电极的CV曲线，利用CV曲线中氢原子脱附峰的电量确定了铂电极的电化学活性面积，考察了不同搅拌速度对氧气饱和的硫酸溶液中铂电极的CV曲线的影响，并根据实验结果绘制了铂阴阳极的直接甲醇燃料电池的输出电压－输出功率曲线。}
\def\keyword{循环伏安法;铂电极;电化学活性面积;直接甲醇燃料电池}

\setCJKmainfont{SimSun}[AutoFakeBold]
\newcommand\xiaobt{\fontsize{14pt}{14pt} \selectfont \noindent}
\newcommand\zhengwen{\fontsize{12pt}{12pt} \selectfont \par}
\newcommand\tupian{\fontsize{10pt}{10pt} \selectfont}
\newcommand\mysection[2]{\xiaobt \bluebf{#1} \vspace*{3pt} \zhengwen #2 \vspace*{14pt} \par}
\newcommand\myfig[2]{\begin{center}
		\includegraphics[width=11cm]{#1} \\
		\tupian #2 \end{center} \zhengwen}

\newcommand\mysi[1]{\si[inter-unit-product = \ensuremath{{}\cdot{}}]{#1}}
\def\my~{\ \textasciitilde}

\begin{document}
	\pagestyle{fancy} \lhead{物理化学实验} \chead{\name} \rhead{\mytoday}
	\begin{titlepage}
		\heiti
\begin{center}
	\includegraphics[scale=1]{logo.eps}\\
	\vspace*{24pt}
	\fontsize{36pt}{0} 物理化学实验报告\\
	\vspace*{72pt}
	\fontsize{24pt}{0}{题目：\underline{\name}}\\
	\vspace*{72pt}
	\fontsize{14pt}{24pt} \selectfont
	姓\qquad 名：\dlmu{\textbf{高靖松}}\\
	学\qquad 号：\dlmu{\textbf{1700011717}}\\
	组\qquad 别：\dlmu{\textbf{第三组}}\\
	实验日期：\dlmu{\textbf{\expday}}\\
	室\qquad 温：\dlmu{\textbf{\exptmp}}\\
	大气压强：\dlmu{\textbf{\expprs}}\\
	\vspace*{64pt}
\end{center}		

\global\mdfdefinestyle{mymdstyle}{%
	leftline=false,rightline=false,
	innerleftmargin=6,innerrightmargin=6,
	backgroundcolor=myyellow,
	linecolor=black,linewidth=0.5pt,%
	leftmargin=0cm,rightmargin=0cm
}
\begin{mdframed}[style=mymdstyle]
\fontsize{12pt}{0} \selectfont
\setlength{\lineskip}{6pt}  \setlength{\parskip}{0pt} 
\noindent \textbf{摘要}\ \ \textnormal{\zhaiyao}\\
\par \noindent
\textbf{关键词}\ \textnormal{\keyword}

\end{mdframed}

	
	\end{titlepage}
	\mysection{$\blacksquare$ 引言}{%研究点
直接甲醇燃料电池（Direct Methanol Fuel Cell，DMFC）是一种质子膜燃料电池，直接以甲醇水溶液或气态甲醇为燃料，具有工作温度低、燃料成分危险性低、结构简单等优点。
%实验方法
循环伏安（Cyclic Voltammetry，CV）法是一种常见的电化学分析方法。该方法通过控制电极电势以不同的速率随时间呈三角波形进行循环扫描，记录扫描过程中的电流-电势曲线，并从中获取丰富的电化学信息。\par
%实验结果
本实验使用循环伏安法研究了甲醇硫酸水溶液中铂电极作为阴阳极的CV曲线，根据实验结果绘制了铂阴阳极的直接甲醇燃料电池的输出电压-输出功率曲线，对改善直接甲醇燃料电池的工作参数提出了指导。 \par }
	\mysection{$\blacksquare$ 实验部分}{\noindent \textbf{1、仪器与试剂} \par
仪器： CHI600B-03099d型电化学分析仪，上海辰华仪器公司；配套软件：CHI600B电化学工作站；三电极系统：参比电极为双盐桥饱和甘汞电极，辅助电极为铂片电极，雷磁公司，工作电极为铂圆盘电极，中国科学院南京土壤研究所\par 
试剂： 0.05 \mysi{mol.L^{-1}} \ce{H2SO4} 溶液； 0.1 \mysi{mol.L^{-1}} \ce{H2SO4} 溶液； 0.05 \mysi{mol.L^{-1}} 甲醇溶液. \par 
\noindent \textbf{2、实验内容} \par 
\noindent \textbf{(1)电极系统的清洁与活化}\par 
依次用自来水和去离子水充分清洗电解池和三电极系统，随后将三电极电解池与 CHI 电化学工作站相连：工作电极连接绿线；辅助电极连接红线；参比电极连接白线，向电解池中加入0.05 \mysi{mol.L^{-1}} \ce{H2SO4} 溶液至液面没过电极底部。\par 
以每秒3个气泡速度向电解池中通入\ce{N2}，并开启磁力搅拌。通气搅拌3分钟至饱和后，在工作站中设置不同的电位区间开始循环扫描，寻找合适的电位窗口（氢区曲线不下沉、氧区曲线不上扬的前提下使区间尽可能宽）。经过多次尝试后得到最佳电位窗口为$-$0.28V\my~ 1.20V，后续活化及测量实验均在该窗口下进行。设置扫描速度为0.5 \mysi{V.s^{-1}}，循环至观察到CV曲线的形状基本不变后停止活化。\par 
CV曲线稳定，且双电层区在0附近，说明电极活化完成。\par 
\noindent \textbf{(2)不同扫描速度下\ce{N2}饱和的\ce{H2SO4} 溶液的CV测试} \par 
活化完成后，停止磁力搅拌和通气体，分别设定扫描速度为0.5、0.2、0.1\mysi{V.s^{-1}}，扫描已用\ce{N2}饱和的0.05 \mysi{mol.L^{-1}}硫酸溶液的CV曲线。 \par 
\noindent \textbf{(3)铂电极表面的氧还原反应（ORR）} \par 
以每秒3个气泡速度向电解池中通入\ce{O2}，并开启磁力搅拌。通气搅拌5分钟至饱和后，分别在高速搅拌并以高流量通入\ce{O2}（每秒5个气泡）、停止搅拌并以低流量通入\ce{O2}（每秒1\my~ 2个气泡）、停止搅拌并关闭气体的条件下，扫描溶液的CV曲线，扫描速度均为0.1 \mysi{V.s^{-1}}。 \par 
\noindent \textbf{(4)铂电极表面的甲醇氧化反应（MOR）}\par 
等体积混合0.1 \mysi{mol.L^{-1}} 甲醇溶液和0.1 \mysi{mol.L^{-1}} \ce{H2SO4} 溶液，加入电解池中。开启搅拌并通入\ce{N2}至饱和后，以0.1 \mysi{V.s^{-1}}的扫描速度对溶液进行循环伏安扫描。\par 
\noindent \textbf{(5)电极系统的清洗}\par 
所有测量实验结束后，重新用自来水和去离子水清洗电极和电解池，在与活化实验相同的条件下再次测量\ce{N2}饱和的0.05 \mysi{mol.L^{-1}} \ce{H2SO4} 溶液的CV曲线。曲线与实验开始时活化后的CV曲线几乎重合，说明电极系统和电解池已经清洗干净。复原实验仪器。 \par}
	\mysection{$\blacksquare$ 数据与结果}{\noindent \textbf{1、扫描速度对CV曲线的影响与电极反应的可逆性} \par
根据实验部分2(2)，绘制了不同扫描速度下的\ce{N2}饱和的0.05 \mysi{mol.L^{-1}} \ce{H2SO4} 溶液的CV曲线，如图1所示（测量10个Segment并取Segment\ 9和Segment\ 10作图）。\par
\myfig{N2_tot_G.eps}{图1\ 不同扫描速度下\ce{N2}饱和的0.05 \mysi{mol.L^{-1}} \ce{H2SO4} 溶液的CV曲线}
从图像中可以看出，随着扫描速度的增加，双电层均在0附近，特征峰位置基本不变，峰电流增大，与循环伏安法的原理相符。在CV曲线中能较明显地找到五个氧化还原峰，分别为电位在0.4 V附近的含氧物种的还原峰，和位于$-$0.2\my~ $-$0.05 V的电位区间的四个H的吸附/脱附峰。扫描速度越大，由于峰电流增大，出峰越明显。根据这些特征峰的峰电流、峰电位值可以研究电极反应的反应性。相关数据列于表1中，其中$i_{p}$为氧区还原峰电流，$E_{p1}$、$E_{p2}$、$E_{p1}'$、$E_{p2}'$分别为氢的强吸附峰、弱吸附峰、弱脱附峰、强脱附峰的峰电位。 \par
\begin{table}[h]
	\centering
	\begin{center}
	\tupian{表1\ 不同扫描速度下\ce{N2}饱和的\ce{H2SO4} 溶液CV曲线的部分参数}
	\end{center}
	\label{tab:width:fixed}
	\sisetup{ 
		 table-number-alignment = center,
		 table-column-width = 2cm , 
	}
	\begin{tabular}
	{
		S
		S[fixed-exponent = -5]
		S
		S
		S
		S
	}
 	\toprule
 	{扫速(\mysi{V.s^{-1}})} & {$i_{p}$ (\num{e-5}A)} & {$E_{p1}$(V)} & {$E_{p2}$(V)} & {$E_{p1}'$(V)} & {$E_{p2}'$(V)} \\ 
 	\midrule 
	0.5&	5.138&	-0.068&	-0.235&	-0.212&	-0.047\\
	0.2&	2.406&	-0.064&	-0.229&	-0.215&	-0.049\\
	0.1&	1.359&	-0.063&	-0.227&	-0.216&	-0.050\\
	\bottomrule
	\end{tabular}
\end{table}
\zhengwen
由表中数据可见，随着扫描速度变慢，对应吸附/脱附峰的差值$\Delta E_{p}$从0.02V减小至0.01V，且峰对称性变高，说明反应基本可逆。另外用氧区还原峰电流对扫描速度做双对数图，具有很好的线性关系，斜率为0.82，接近于1，说明此条件下的含氧物种还原峰具有可逆吸附波的特征。\par

\noindent \textbf{2、铂电极电化学活性面积的测定} \par
在0.1 \mysi{V.s^{-1}}的扫描速度下，停止搅拌且关闭气体，绘制被\ce{N2}饱和的硫酸溶液的CV曲线，如图3所示。根据氢脱附区的曲线下围面积（以右侧双电层的最低点为基线），用Origin给出其积分面积为\num{8.63e-7} \mysi{V.A}。根据电量$Q$的定义，有
$$Q=\int{I\mathrm{d} t}=\int{\frac{I}{v}\mathrm{d} E}=\frac{1}{v}\int{I\mathrm{d}E}=\frac{Area}{v}$$\par 
因此$Q=\num{8.63e-6}$C，多晶Pt表面满单层H原子脱附量为0.21 \mysi{mC.cm^{-2}}，故Pt电极的电化学活性面积为0.041 \mysi{cm^{2}}。\par 
\myfig{DesorpArea_G.eps}{图2\ 被\ce{N2}饱和的硫酸溶液的CV曲线中氢脱附区的积分面积图}

\noindent \textbf{3、氧还原反应（ORR）的CV曲线} \par
在0.1 \mysi{V.s^{-1}}的扫描速度下，停止搅拌且关闭气体，分别绘制被\ce{N2}和\ce{O2}饱和的硫酸溶液的CV曲线，并求出差值，结果如图3所示（测量8个Segment并取Segment\ 7和Segment\ 8作图）。\par
由图可见，当电位大于0.6V时，介质中的\ce{O2}几乎不发生额外的化学行为，CV曲线与\ce{N2}的基本重合。当电位回扫到0.4V时，还原峰电流迅速增大，此时电极上发生了ORR反应，以差值电流超过\num{2.0e-7}A为界限，ORR发生的起始电位为0.532V。\par 
控制扫描速度为0.1\mysi{V.s^{-1}}，分别在高速搅拌并以高流量通入\ce{O2}（每秒5个气泡）、停止搅拌并以低流量通入\ce{O2}（每秒1\my~ 2个气泡）、停止搅拌并关闭气体的条件下，测定溶液的CV曲线，如图4所示。\par 
同样在电位大于0.6V时，三条曲线几乎重合，说明此时溶液的电化学响应基本不变。但在发生ORR的部分，由于溶液中对流、传质情况的不同而产生的较大差别。对流越剧烈，电极附近的\ce{O2}和\ce{H^{+}}浓度涨落就越大，使得电流变得不稳定。\par 
\myfig{N2O2diff_G.eps}{图3\ 被\ce{N2}和\ce{O2}饱和的0.05 \mysi{mol.L^{-1}} \ce{H2SO4} 溶液的CV曲线及差值曲线}
\myfig{O2_tot_G.eps}{图4\ 不同搅拌速度下被\ce{O2}饱和的0.05 \mysi{mol.L^{-1}} \ce{H2SO4} 溶液的CV曲线}

\noindent \textbf{4、甲醇氧化反应（MOR）的CV曲线} \par
在0.1 \mysi{V.s^{-1}}的扫描速度下，停止搅拌且关闭气体，分别绘制被\ce{N2}饱和的硫酸溶液和甲醇/硫酸混合溶液的CV曲线，并求出差值，结果如图5所示（测量8个Segment并取Segment\ 7和Segment\ 8作图）。\par
被\ce{N2}饱和的甲醇/硫酸溶液的CV曲线较为复杂，与原始的\ce{N2}饱和的硫酸溶液的CV曲线相差较大，正扫时存在两个甲醇氧化峰，分别位于0.5V和1.2V右侧，回扫时除ORR还原峰外，还存在一个0.3V左右的氧化峰，意味着在回扫阶段依然是甲醇的氧化反应而不是还原反应。取阳极区段中差值电流为0的点，可以读出MOR发生的起始电位为0.084V。
\myfig{MeOHN2diff_G.eps}{图5\ 被\ce{N2}饱和的硫酸溶液和甲醇/硫酸混合溶液的CV曲线及差值曲线}

\noindent \textbf{5、铂阴阳极的直接甲醇燃料电池的输出功率-输出电压曲线} \par
将铂电极上发生的ORR和MOR组合成原电池。在等电流的情况下，输出电压为ORR电位减去MOR电位，输出功率为输出电压乘以电流，即
$$U_{cell}=E_{ORR}-E_{MOR},P_{cell}=i\cdot U_{cell}$$\par 
分别将扣除\ce{N2}背景后的氧气和甲醇的CV曲线中ORR和MOR的区段绘制在图中，结果如图6所示。由于电池两极中电流的方向相反，因此需要将ORR曲线向上翻转，即电流取绝对值。为保证电池能够正常输出电流，氧气应在正极还原，甲醇应在负极氧化，输出电压为正值，因此只有图中的阴影区域才是这个电池的有效工作区域。\par
\myfig{EffArea_G.eps}{图6\ 铂阴阳极的直接甲醇燃料电池的有效工作区域}
首先提取出阴影区域的电位电流信息，交换$xy$轴后，并将曲线分成上下两部分。用Origin的插值功能将原本不均匀分布的$x$值填充为等间隔的值，以方便后续求差值。用两组曲线上插值后的点两两相减，求出对应的输出电压$U_{cell}$和输出功率$P_{cell}$。作出电池的输出功率-输出电压曲线，如图7所示。\par
\myfig{Power_G.eps}{图6\ 铂阴阳极的直接甲醇燃料电池的有效工作区域}
从图中可以看出，铂阴阳极的直接甲醇燃料电池的输出功率随着输出电压的升高，先升高后降低，在输出电压约为0.145V时达到最大值\num{8.17e-7} W。\par

}
	\mysection{$\blacksquare$ 结果与讨论}{\noindent \textbf{1、误差分析和实验的改进} \par
在本次实验中，所有的数据和结论均是由循环伏安曲线图直接或间接的得出，因此误差的来源即为循环伏安测试的过程。在使用电化学工作站进行循环伏安测试时，电流的灵敏度参数的选取很重要，当电解质溶液的导电性较好时，为防止电流过大超过量程，需要选取较大的灵敏度参数，即更低的灵敏度；而如果导电性较差电流很小时，就需要选取较小的灵敏度参数，提高仪器检测电流的灵敏度，否则可能出现实验数据波动较大，精度较差等问题。\par 
\noindent \textbf{2、本次实验中的重大失误} \par
在本次实验中，由于预习不够充分，没能准确理解实验讲义中“在实验过程中参比电极的二次盐桥中加 入的始终是 0.05 M 硫酸水溶液；实验结束后更换为去离子水”的含义。在实验开始时没有将二次盐桥中的去离子水替换为硫酸，仅在测试结束后换成了去离子水。\par 
从参比电极的结构角度考虑，如果二次盐桥内放去离子水，且长时间放置，会使得内盐桥中的一部分氯化钾扩散至二次盐桥中，这会在一定程度上影响参比电极的电导，可能使得其上有电流流过，造成工作电极和辅助电极间电位、电流的测量结果不准确。因而实验中的CV图均不够准确，例如双电层略微向下偏离0水平，氢区存在双线电流的情况。\par 
今后使用有二次盐桥的电极时，应注意及时将盐桥内的电解液换成与电解池中溶质种类、浓度均相同的溶液。\par
\noindent \textbf{3、实验结论} \par
本实验用循环伏安法研究了不同扫描速度下\ce{N2}饱和的0.05 \mysi{mol.L^{-1}} \ce{H2SO4} 溶液中铂电极的电化学响应性质，证明了含氧物种还原峰是可逆吸附波，铂电极的电化学活性面积为0.041 \mysi{cm^{-2}}。测定了\ce{N2}饱和的甲醇/硫酸混合溶液以及不同搅拌速度下\ce{O2}饱和的0.05 \mysi{mol.L^{-1}} \ce{H2SO4} 溶液的CV曲线，得出ORR发生的起始电位为0.532V，MOR发生的起始电位为0.084V，证明在对流剧烈的情况下，ORR反应的电流不确定性显著增大。结合MOR和ORR的CV曲线，绘制出铂阴阳极的直接甲醇燃料电池的输出功率-输出电压曲线，得到输出电压约为0.145V时，电池输出功率达到最大值\num{8.17e-7} W。\par }
	\mysection{$\blacksquare$ 参考文献}{\noindent [1]\ 北京大学化学学院物理化学实验教研组，物理化学实验(第4版)，北京大学出版社，\textbf{2002}.\par
\noindent [2]\ 贾梦秋，杨文胜. 应用电化学，高等教育出版社，\textbf{2004}.\par
\noindent [3]\ 胡会利，李宁主编，电化学测量，国防工业出版社，\textbf{2007}.\par}
	\clearpage
\end{document}